\section{Power-Up and Test Precautions}
\label{sec:method-safety}

Bring-up used a controlled, current-limited source, rated probes, separation
from exposed conductors, and an external input disconnect. Hardware, source,
load, and probe references were checked with the burst-enable output PA10
inhibited; UART and oscilloscope connections were completed before energizing.
Probe differential and common-mode ratings were observed, and grounded
passive-probe clips were excluded from line-referenced nodes. Source isolation
did not remove shared-reference constraints.

Shutdown removed input power and the enable command. Residual DC-link and
output-capacitor voltages were checked before handling. PA10 inhibition stops
commanded transfer but does not discharge stored energy. The external
disconnect provided an emergency stop independent of UART communication.

\begin{table}[htbp]
\centering
\caption{Connections and settings needed for powered tests.}
\label{tab:method-safety-metadata}
\begin{tabularx}{\textwidth}{@{}p{0.25\textwidth}X@{}}
\toprule
Item & Settings and connections\\
\midrule
Input source & Voltage, frequency, current limit, and isolation/reference arrangement\\
Probe connection & Node pair, attenuation/range, and voltage/common-mode rating\\
Output load & Load mode, setpoint, selected output setting, and connection state\\
Shutdown & External disconnect and PA10 inhibit state\\
Stored energy & Discharge interval and verified residual voltage at accessible nodes\\
\bottomrule
\end{tabularx}
\end{table}
